\documentclass[12pt]{article}

\usepackage[a4paper, total={6in, 8in}]{geometry}
\usepackage{textcomp}

\begin{document}
\setlength{\parindent}{0pt}

\def \t {\textrightarrow}
\def \v {\vspace{1em}}
\def \bi {\begin{itemize}}
\def \ei {\end{itemize}}
\def \s[#1] {\section*{#1}}
\def \ss[#1] {\subsection*{#1}}
\def \sss[#1] {\subsubsection*{#1}}

\s[Tacito]
Impero, le critiche che arrivano per maldicenze (tipico dell'impero) \\
Vengono mosse accuse quando pubblica l'Agricola \t ma anche quando si dedica agli annales \t lui è lo storico che ha a cuore un metodo di indagine preciso, nitido \\
Ha anche interessi per la germania di tacito, intesa come opera etnogeografica \t osservazione diretta dei fenomeni, ovvero nel caso dell'oratoria analizza l'oratore \t che non è piu il tempo della oratoria ciceroniana

\ss[Dialugs de oratoribus]
Tacito parla dell'importanza formativa dell'oratoria \t oratoria forma la classe dirigente, e nel "Dialogus de oratoribus" parla di fiamma \t ma manca nell'impero il pluralismo delle opinioni \\
"Dialugs de or." \t è tramandato nella tradizione scritta insieme all'agricola e alla germania \t ha raccolto in unico manoscritto tutte e tre queste opere \\
Fino pero al 16 secoloco, i filologi mettono in dubbio l'autenticita del dialogus (forse non l ha composto lui) \t qua infatti si riflette un modello neociceroniano (forbito ma non prolisso) \t scuola che devia dal cicerone originario, ma che ne mantiene alcuni aspetti \t la rigorosita severa dell'armonia viene un po trascurata, ma non ci sono tante variationes \\
Tacito invece è pieno di variationes \t c'è scollamento \\
Garbarino \t insieme a Gian Biagio Conti, insieme a Pontigia \t alti nomi della filologia latina \\
Altri pensano che il prodotto sia di tacito, ma che ci siano delle disarmonie nella datazione \t il Dialogus sarabbe stato scritto in eta giovanile, e quindi influenzato dalla scuola ciceroniana \\
Dialougs risale al 75/77 \t sappiano il contesto in cui il dialogus avviene, e si svolge tra: Curiatio, Marco Apro, Vipstano Messalla e Giulio II \\
L'oggetto del contendere è che si da + importanza alla poesia che alla oratoria \t la poesia è + libera, mentre oratoria è + vincolata \\
Dopo un dibattito che riguarda la poesia, la decadenza dell'eloquenza \t l'autore fa una scorsa meno pregnante di quanto sia veloce sul deterioramento dell'educazione \\
Una cultura che sia mirata (ed. singolare) viene sostituita da una educazione lacunosa generalizzata \\
È un opera di oratoria che non trascura l'oggetto del dire \t ma la riflessione si fa poi amara \t la consapevolezza di tacito ha una parola: anarchia \\
L'oratoria nel suo tempo scade in un anarchia, impensabile al tempo di cicerone \\
Machiavelli \t si parla di tacitismo \t tacito viene accusato di dove???????? \t tacito parla di un oratoria destinata a spegnersi a causa della prevaricazione del potere \\
Roma ha perso la sua assertivita \t rimane il tacitismo = aderire senza distinguersi, che è la paura + grande di tacito \\
Inoltre il principato non puo essere neutralizzato \t secondo la visione di Materno, una soffocata voce di un urlato impero \t l'oratore non puo fare nulla, ha campo d'azione ristretto e non puo far sentire la sua voce \\
Manca il contrasto diretto, e il terreno dove il contrasto è costruttivo e non offensivo \t uno stato come lui lo intende non puo essere applicabile, ma questo l'aveva gia detto nell Agricola

\ss[L'Agricola]
Agricola è un'opera biografica encomiastica celebrativa, ma anche la volonta di affermare un modello umano di comportamento \t un modello assertivo di uomo politico e distintivo di classe dirigente \\
Siamo nell eta di Traiano \t il tentativo di una boccata d'aria dopo i il periodo di Domiziano, e l'attenzione di tacito verso il suocero \\
Laudatio funebris \t lode riferito a Giulio Agricola, che è suocero che gli da in sposa la figlia \\
Domiziano e traiano si succedono \t lui scrive l opera dopo aver vissuto sul campo \\
Da ritratto di agricola nelle sue potenzialita \t nella sua capacita di servire lo stato con fedelta e nel contraddistinguersi per virtus \\
Nel periodo domizianeo in particolare, quando il terrore imperversava \t ma agricola è rimasto incorrotto durante la corruzione, e di morire silenziosamente \t agricola muore e non si sanno le cause della sua morte \t alcuni dicono che è stato ucciso da Domiziano \\
Apologia della potenzialita costruttiva di stare nella classe dirigente \t e di ampliare i confini di roma \\
Il carattere composito dell'opera è imp. \t laudatio funebris ma non tanto, opera sprona all'emulazione \t materiali storici integrati a materiali biografici, aneddoti, composizione che mira alla descrizione delle sue caratteristiche senza enfasi eccessiva e senza mistificazione \\
Un modello ideale? \t si ma radicato sul reale \t questo fa capire che ci sono tratti sallustiani (x variatio credo) e un influenza ancora ciceronia, affettuosita che si vede (vuole veramente bene ad agricola) \\
Esempio di un uomo politico in un epoca di corruzione, e quindi una presenza di eccezionalita che ricorda lo stile di Livio \t un'opera composita in cui i modelli precedenti \\
Valore educativo dell'opera (e anche l'opera di cicerone educava, movere docere flectere), ed esalta l uomo per le sue qualita \t e si ha la passione per il ritratto, che ha dietro sallustio \\
Stile difficile accumana tacito a stallustio, che tralaltro sono anche entrambi storici, e anche Livio \\
\textbf{I tre metodi di storia:}
\bi
  \item sallustio: monografia storica
  \item livio: miracula portenta et prodigia \t parla della storia in modo gonfiato, parla di fenomeni eccezionali
  \item tacito: metodo storico diretto \t sempre fedele a un indagine storica diretta, interagisce con i fatti \t e lo fa con l'osservazione diretta e con l'analisi delle fonti
\ei
Il filologo fa come il restauratore \\
Tacito ha amore per il vero \t ed analizza la storia \\
In Cicerone "Della proelio"?? \t anche qua cicerone ha amore per il vero ??
\end{document}
